\documentclass[10pt]{article}
\usepackage[margin=20mm]{geometry}
\usepackage{amsmath,amssymb,mathtools,fontspec,unicode-math,xurl,hyperref}
\IfFontExistsTF{DejaVu Serif}{\setmainfont{DejaVu Serif}}{\setmainfont{Latin Modern Roman}}
\setmathfont{latinmodern-math.otf}
\setlength{\parindent}{0pt}\setlength{\parskip}{6pt}\sloppy
\title{Technical Verification of
\(\mathbb H^{p,q}\)-Convex Cocompactness and
Higher Higher Teichmüller Spaces}
\author{AI-assisted technical review}
\date{}
\begin{document}
\maketitle
Abstract: This report completes a bounded recovery of an unfinished initial review of Beyrer and Kassel’s arXiv version 7. An explicit product-group example disproves the positivity extension in Lemma 6.7 and the deletion rule in Corollary 6.8(2). The deletion rule is used in the proof of the principal non-degeneracy proposition, so a replacement justification is required. This is a Major proof-level finding; no counterexample to a main theorem is claimed. Two further local overstatements admit restricted repairs. A fourth finding concerns the literal real tensor-product assertion in Fact 6.9; its additional semiproximal application is not refuted. A boundary-dimension ambiguity is recorded separately. Scientific adjudication and independent acceptance remain pending.\par
\section*{Identity, source and scope}
The paper is Jonas Beyrer and Fanny Kassel, \(\mathbb H^{p,q}\)-Convex Cocompactness and Higher Higher Teichmüller Spaces, associated with the frozen 2025 journal record Geometric and Functional Analysis, DOI 10.1007/s00039-025-00719-9 (https://doi.org/10.1007/s00039-025-00719-9). The reviewed source is arXiv:2305.15031v7 (https://arxiv.org/abs/2305.15031v7), submitted 3 October 2025, with internal date 6 October 2025. It has 76 physical PDF pages; the printed page numbers used below agree with physical pages. No equivalence with the journal version of record is assumed.\par
The earlier checkpoint records full text reading; that historical reading is not represented here as a completed proof audit. Recovery concentrated on its open candidates: the boundary extension in Section 3, the crown orbit argument in Section 5, the representation theory and non-degeneracy argument in Section 6, and the Appendix A boundary case. The source definitions, operative hypotheses, decisive pages and downstream dependencies of these claims were checked. Selected statements and proofs in Benoist’s retained author manuscript [Benoist] were consulted.\par
This is an initial technical report with a declared scope, not a certification of every theorem, every cited dependency, or the independent acceptance of the paper. The coverage ledger distinguishes historical reading from the present focused checks. No author program was executed. An unavailable separate code release would not by itself constitute a mathematical defect.\par
\section*{The reducible positivity extension and deletion rule}
Problem 1 (An invalid deletion rule used in the non-degeneracy proof).  Lemma 6.7 (p. 52) and Corollary 6.8(2) (p. 53) are false as stated for reducible representations. Their use in Proposition 6.12 and the proof of Proposition 1.13 creates a Major proof-level issue.\par
\subsection*{Original hypotheses}
In Definition 6.3, a representation \(\tau\) of the real points \(G\) of a connected reductive real algebraic group is positively \(\Gamma\)-semi-proximal when, for every \(\gamma\in\Gamma\) loxodromic in \(G\), the matrix \(\tau(\gamma)\) has a positive real eigenvalue of maximal modulus. Simplicity of that eigenvalue is not required. Loxodromic means that the Jordan projection belongs to the interior of the positive Weyl chamber.\par
Lemma 6.7 assumes that \(\Gamma\) is a Zariski-dense subgroup, that \(G'\subset G\) is an open subsemigroup meeting \(\Gamma\), and that \(\tau\) is positively \((\Gamma\cap G')\)-semi-proximal. It concludes global positive \(\Gamma\)-semi-proximality. There is no irreducibility hypothesis. Corollary 6.8 assumes irreducible summands \(\tau_i\) with highest weights \(\chi_i\) and a positively \(\Gamma\)-semi-proximal direct sum. Part (2) says that if some \(X\) in the limit cone satisfies \(\chi_1(X)>\chi_2(X)\), one may delete \(\tau_2\) and retain positive \(\Gamma\)-semi-proximality.\par
\subsection*{A counterexample satisfying all of these conditions}
Equation reference: eq:group, eq:reps. Let \[
 G=\mathrm{SL}_2(\mathbb R)\times\mathrm{SL}_2(\mathbb R),\qquad
 \Gamma=\mathrm{SL}_2(\mathbb Z)\times\mathrm{SL}_2(\mathbb Z),\] and, with the subscript identifying the factor, set \[
 \tau_1=\mathrm{Sym}^2\mathrm{std}_1,\qquad
 \tau_2=\mathrm{Sym}^2\mathrm{std}_2,\qquad \tau_3=\mathrm{std}_2.\] These are irreducible real algebraic representations of the product, of dimensions \(3,3,2\). The algebraic group is connected reductive. The subgroup is discrete and Zariski dense: the integer upper and lower unipotent subgroups are Zariski dense in the corresponding one-parameter root subgroups, which generate each \(\mathrm{SL}_2\) factor.\par
For every \(B\in\mathrm{SL}_2(\mathbb R)\), the matrix \(\mathrm{Sym}^2 B\) has a positive eigenvalue \(\rho(B)^2\), where \(\rho\) is spectral radius, and \(\rho(B)\geq1\). Indeed, for real reciprocal eigenvalues the squares supply this eigenvalue; for a conjugate unit-modulus pair the product eigenvalue is \(1\). This also covers repeated eigenvalues. Consequently, \(\mathrm{Sym}^2 B\) has spectral radius at least that of \(B\). For any \((U,V)\in G\), both symmetric-square blocks have a positive maximal-modulus eigenvalue and the second one dominates \(\mathrm{std}_2(V)\). Thus \(\tau_1\oplus\tau_2\oplus\tau_3\) is positively semiproximal on all of \(G\), which is stronger than the corollary’s hypothesis.\par
Take \[A=\begin{pmatrix}2&1\\1&1\end{pmatrix},\qquad
 r=\frac{3+\sqrt5}{2}>1.\] The eigenvalues of \(A\) are \(r,r^{-1}\). In Weyl chamber coordinates \((x,y)\) the weights in [eq:reps] (\#eq:reps) are \(2x,2y,y\). The element \[X=\lambda(A^2,A)=(2\log r,\log r)\] belongs to the limit cone, and \(\chi_1(X)=4\log r>2\log r=\chi_2(X)\). All the assumptions of Corollary 6.8(2) therefore hold.\par
Equation reference: eq:spectrum. Now evaluate the proposed surviving representation at \(\gamma=(A,-A^3)\in\Gamma\). Both factors are hyperbolic, so \(\gamma\) is loxodromic in \(G\). Its spectrum on \(\tau_1\oplus\tau_3\) is \[
 r^2,\;1,\;r^{-2},\;-r^3,\;-r^{-3}.\] Since \(r^3>r^2\), the unique maximal-modulus eigenvalue is negative. The deleted block instead has positive leading eigenvalue \(r^6\). This proves the claimed discrepancy without a numerical eigenvalue test. The preserved exact computation additionally verifies \[A^2=\begin{pmatrix}5&3\\3&2\end{pmatrix},\qquad
 -A^3=\begin{pmatrix}-13&-8\\-8&-5\end{pmatrix}\] and the polynomial identities for \(r,r^2,-r^3\).\par
Equation reference: eq:semigroup. For completeness, the same example directly refutes Lemma 6.7. Use \(\tau=\tau_1\oplus\tau_3\). For a positive matrix \(U\), let \(m(U)\) denote its minimum row sum, and define \[
 G'=\{(U,V)\in G:\ U_{ij}>1\ \hbox{for all }i,j,\
                  \|V\|_\infty<m(U)^2\}.\] This set is open. It meets \(\Gamma\) at \((A^2,A^2)\), since \(m(A^2)=5\) and \(\|A^2\|_\infty=8<25\). If \((U,V),(U',V')\in G'\), every entry of \(UU'\) exceeds \(2\), and \[m(UU')\geq m(U)m(U'),\qquad
 \|VV'\|_\infty\leq\|V\|_\infty\|V'\|_\infty
       <m(U)^2m(U')^2\leq m(UU')^2.\] Thus \(G'\) is a subsemigroup. Perron–Frobenius theory gives \(\rho(U)\geq m(U)>0\) and a positive leading eigenvalue of \(U\). Therefore the positive eigenvalue \(\rho(U)^2\) of \(\mathrm{Sym}^2U\) strictly dominates \(\rho(V)\) on \(G'\). The representation \(\tau\) is positively semiproximal there, yet [eq:spectrum] (\#eq:spectrum) disproves its global positive \(\Gamma\)-semi-proximality.\par
\subsection*{The precise proof failure and the valid irreducible case}
The proof of Lemma 6.7 fixes the highest eigenspace \(V^+\) of \(\tau(a)\) and then, using [Retained latex: \textbackslash{}cite[Corollary 8.5]\{Benoist\}], replaces \(ta\) by \(t'a'\) with \(t'\) close to \(t\). For a reducible representation, changing \(a\) to \(a'\) can change which constituent supplies the highest eigenspace. The fixed \(V^+\) need not be the highest eigenspace of \(\tau(a')\). The determinant condition on that old space consequently does not imply the asserted spectral conclusion.\par
This objection does not refute the irreducible sign argument in [Retained latex: \textbackslash{}cite[Lemma 6.3]\{Benoist\}], consulted on pp. 27–28, or the exposed single-summand conclusion of Corollary 6.8(1). The source’s part (1) has an independent irreducible justification; it should not be discarded merely because the new general lemma is false. The compact-factor conclusion of Benoist’s Corollary 8.5, consulted on p. 46, does not control a changing dominant constituent.\par
\subsection*{Core dependency, alternative readings and repair limits}
Proposition 6.12, Step 1 (p. 55), explicitly invokes Corollary 6.8(2) to delete one member of a pair of non-isomorphic dual constituents under a minimality argument. Step 2 invokes it again on p. 56. In the proof of Proposition 1.13 (p. 58), a minimal subrepresentation of \(\kappa\oplus\kappa^*\) is selected, and the text invokes the argument of Step 1 to obtain disjoint index sets, the dimension bound \(\dim\tau\leq k\), and then the contradiction \(p\leq\delta(\tau)+p-k<p\). The introduction identifies Proposition 1.13 as the second main ingredient of Theorem 1.3 (p. 9); p. 10 explicitly uses it in Theorem 1.17 and then in the closedness deduction for Theorem 1.3. The invalid rule is therefore tied to a central derivation, rather than an unused general observation.\par
There is a significant limit to this conclusion. The counterexample above does not instantiate the special dual-pair and minimality configuration of Step 1. It disproves the rule cited to justify that step, not the specialized step under every additional hypothesis. It is not a counterexample to Proposition 6.12, Proposition 1.13, or a main theorem. A specialized replacement lemma might restore the argument, but none is proved in this report.\par
One local improvement is available. In Step 2, once the retained exposed constituents have individually been shown positive by Corollary 6.8(1), their direct sum is positive by Remark 6.5(2). The cone-preserving representation \(\upsilon\), when nontrivial, is also positive by Fact 6.11(2). This replaces the positivity sentence on p. 56. It does not by itself justify Step 1, the remaining minimality and kernel arguments, or the full proposition. The supported assessment is Major for the demonstrated invalid core-related derivation, with a complete repair unverified.\par
\section*{Strict contraction at the ideal boundary}
Problem 2 (A strict interior contraction need not remain strict on the boundary).  The proof of Proposition 3.28(2), p. 33, overstates the strictness of a continuous boundary extension. This is a local Minor finding; the proposition has the following restricted repair.\par
Here strict \(1\)-Lipschitz means a strict distance inequality for every pair of distinct points, as on p. 28; it does not mean a single uniform Lipschitz constant less than \(1\). Use \(p=2,q=1\) and the map from Example 3.23(ii), p. 29, on the open upper unit hemisphere: \[d=y_1^2-y_2^2,\qquad
 f(y_0,y_1,y_2)=
 \left(\sqrt{\frac{1+d}{2}},\sqrt{\frac{1-d}{2}}\right)\in S^1.\] It is equivalently the source’s formula \((\sqrt{y_0^2/2+y_1^2},\sqrt{y_0^2/2+y_2^2})\). The angular coordinate \(\theta=\tfrac12\arccos d\) satisfies \[\|\nabla_{S^2}\theta\|^2
   =\frac{1-y_0^2-d^2}{1-d^2}
   =1-\frac{y_0^2}{1-d^2}<1 .\] For two interior points the minimizing spherical geodesic lies in a compact subset of the open hemisphere, so integration gives strict contraction. This verifies the operative interior assumption. On the boundary arc \(y(t)=(0,\cos t,\sin t)\), \(0<t<\pi/2\), however, \(f(y(t))=(\cos t,\sin t)\). Distances on this arc are preserved. The continuous extension is thus not strictly \(1\)-Lipschitz on all boundary-boundary pairs.\par
The proof only needs strictness between an interior point \(x\) and a boundary point \(z\). Choose an interior point \(y\) strictly between them on a minimizing geodesic. The extension is nonexpansive, and \[d(f(x),f(z))
 \leq d(f(x),f(y))+d(f(y),f(z))
 <d(x,y)+d(y,z)=d(x,z).\] This proves precisely the needed replacement and preserves the conclusion of Proposition 3.28(2). No failure of that proposition or of a central theorem follows from the overstatement.\par
\section*{Straightness of crown orbits}
Problem 3 (Metric geodesicity does not imply projective straightness).  The proof of Proposition 5.9, p. 45, asserts projective straightness for every nonzero Cartan direction. The assertion is false in that generality, but the directions used in the quadrilateral construction satisfy a valid stronger restriction. The supported impact is Minor.\par
Use a \(2\)-crown in signature \((2,2)\), corresponding to \(p=2,q=1\), with its isotropic crown basis and positive orthant \(\mathcal O\). Take \(x=[1:1:1:1]\in\mathcal O\) and \(a=(1,0)\). In the diagonal coordinates of (5.1), p. 43, the orbit is \[x(t)=[e^t:1:e^{-t}:1].\] At \(t=0,\log2,\log4\), the first three coordinates of representatives have determinant \[\det\begin{pmatrix}1&1&1\\2&1&1/2\\4&1&1/4\end{pmatrix}
 =-\frac34\ne0.\] Their vector span has dimension at least three; the three projective points cannot lie in a projective line. All representatives remain in the positive orthant and the action is the stated crown action. The source’s Hilbert-distance formula (5.3) still gives a metric geodesic. It is the additional projective-line assertion that fails.\par
For the construction actually used on p. 45, the directions have coordinates of equal absolute value. In this case \(\exp(ta)x=e^{ct}u_++e^{-ct}u_-\) for two fixed vectors, and the orbit does lie in their projective span. Restricting the straightness statement to these directions suffices for that construction. This evidence does not refute Proposition 5.9 or its use in Proposition 5.10 on p. 46.\par
\section*{Real tensor products and a quaternionic case}
Problem 4 (The literal factorization in Fact 6.9 omits a real quaternionic case).  Fact 6.9, p. 53, states an \(\mathbb R\)- or \(\mathbb C\)-tensor factorization for every irreducible real representation of a product of connected reductive real algebraic groups. Its generality is false. The present example does not refute the additional semiproximal case used later; this is a Minor general-statement finding.\par
Let \(G=\mathrm{Sp}(1)\times\mathrm{Sp}(1)\) be the real algebraic product of unit quaternion groups, and let \(V=\mathbb H\) as a four-dimensional real vector space, with action \[(a,b)\cdot x=axb^{-1}.\] The action is real algebraic, and the groups are connected and reductive. The representation remains irreducible when restricted to either factor: the orbit of any nonzero quaternion under left (or right) multiplication spans \(\mathbb H\) over \(\mathbb R\).\par
Suppose the factorization stated in Fact 6.9 existed: \(V=W_1\otimes_{\mathbb K}W_2\), \(\mathbb K=\mathbb R\) or \(\mathbb C\), with each \(W_s\) irreducible as a real representation of its factor. On restriction to the first factor, this tensor product is a direct sum of copies of the underlying real representation \(W_1\); the same statement holds for \(W_2\) and the second factor. The two irreducible restrictions of \(V\) would force each \(W_s\) to be isomorphic to a four-dimensional real representation. But the proposed tensor product then has real dimension \(16\) over \(\mathbb R\), or \(8\) over \(\mathbb C\), rather than \(4\). This is the missing quaternionic-type case.\par
The retained version of [Retained latex: \textbackslash{}cite[Lemma 5.5]\{Benoist\}], p. 22, contains the corresponding general assertion; citing it does not remove this explicit dimensional obstruction. This report makes no first-discovery claim. A corrected real representation-theoretic statement requires the relevant quaternionic/real-form treatment, or a proved restriction to an adequate special class.\par
There is also decisive counterevidence against overstating the application. The compact example is not \(G\)-semiproximal: left multiplication by a nonreal unit quaternion, with the second factor the identity, can have no real eigenvalues. Proposition 6.12, Step 3, has additional semiproximal hypotheses. The example therefore disproves the unrestricted Fact 6.9 but establishes no failure of that special case or central conclusion. No complete replacement theorem is claimed.\par
\section*{Appendix A scope clarification and overall assessment}
Lemma A.8(3), pp. 68–69, uses an inclusion \(\mathbb R^{2,1}\subset\mathbb R^{2,q}\). If equality is permitted, the written \(q\geq1\) scope admits \(q=1,c=0\) and, for the form \(\operatorname{diag}(1,1,-1)\), \[v_0=(0,5/4,3/4),\qquad v_2=(0,5/4,-3/4),\qquad v_1=(1,0,0).\] All three vectors have norm \(1\), the two required pairings with \(v_1\) are zero, and the spanning determinant has absolute value \(15/8\). The pointwise stabilizer of \(v_0,v_2\) in \(\mathrm{SO}(2,1)\) is the identity: it fixes their nondegenerate two-dimensional span, and determinant one forces identity on the remaining line. It cannot move \(v_1\) out of \(\mathbb R^{2,0}\) as requested. A proper-inclusion interpretation excludes this example, and the actual application on p. 71 already separates \(q=1\) and works with \(q\geq2\). This is recorded as a scope clarification, not an unconditional defect in that application. A prior \(c=0\) concern about Lemma A.9 failed its premises and is not promoted as a finding.\par
The highest supported assessment is Major, attached to Problem 1 and its explicit role in the non-degeneracy/closedness proof. Problems 2 and 3 have proved local repairs. Problem 4 requires qualification but supplies no counterexample under its later semiproximal application. No result here establishes that the paper’s principal theorems are false. No author adoption or independent acceptance is asserted.\par
Item | Supported assessment | Repair state\par
Positivity deletion | Major; false general rule, core proof dependency | Complete core repair unverified\par
Boundary strictness | Minor local overstatement | Interior-boundary repair proved\par
Orbit straightness | Minor local overstatement | Needed restricted directions proved\par
Real tensor factors | Minor general assertion | Special-case repair unverified\par
Appendix A dimension | Scope clarification | Actual application excludes \(q=1\)\par
Scientific decision is pending. Human reviews are empty. An independent scientific acceptance must be performed later by a different authorized scientific conversation.\par
9 J. Beyrer and F. Kassel, \(\mathbb H^{p,q}\)-Convex Cocompactness and Higher Higher Teichmüller Spaces, arXiv:2305.15031v7 (3 October 2025), https://arxiv.org/abs/2305.15031v7. Journal identity DOI:10.1007/s00039-025-00719-9; version equivalence unverified.\par
Y. Benoist, Convexes divisibles III, retained author manuscript, 48-page PDF, https://www.imo.universite-paris-saclay.fr/\textasciitilde{}yves.benoist/prepubli/01divisibleIII.pdf. Consulted specifically: Lemma 5.5 (p. 22), Lemma 6.3 (pp. 27–28), and Corollary 8.5 (p. 46).\par
\section*{Supporting checks}
> .\par
The retained Benoist PDF SHA256 is\par
> .\par
The recovery directory adds the independent exact-rational/quadratic-field checks, actual source-image view receipts, coverage, findings, build payload and execution records. The exact computations support the concrete matrix and determinant identities; the universal arguments in this report are analytical proofs and are not inferred from finite tests.\par
All current scientific reasoning, report writing, calculations, compilation and actual page inspection were performed in the current scientific conversation. The controller-supplied UI model observation and the runtime’s self-description are recorded separately in provenance; no independent verification of their backend equivalence is claimed. The retained earlier work remains attributed to the original conversation. No human responsibility declaration or institutional endorsement is supplied.\par
\end{document}
